\section{Test Objectives and Configuration}
\label{sec:method-objectives}

The test campaign assessed loaded regulation, startup, load removal, and
no-load energy control. Input power quality, efficiency, switching waveforms,
and surface temperatures were measured to identify operating limits and
priorities for the next revision.

The V1.0 PCB was tested with two series resonant inductors. V1.1 is the corrected
schematic reference. The baseline voltage-control firmware is described in Section~\ref{sec:ctrl-firmware-identities}. Where available,
UART \texttt{VER}, \texttt{STATUS}, and \texttt{EXPORT} output identify the
active version, mode, calibration, and thresholds for comparisons of burst
operation.

\begin{table}[htbp]
\centering
\caption{Prototype tests and measurement objectives.}
\label{tab:method-objectives}
\begin{tabularx}{\textwidth}{@{}p{0.31\textwidth}X@{}}
\toprule
Question & Measured quantities\\
\midrule
Output capability & VBUS, load current, and delivered output power\\
Control response & Output voltage and PA10 during startup, steady operation, and load removal\\
Input-current shaping & Input voltage/current, real power, PF, THD, and harmonics\\
Conversion loss & Input and output power at the same recorded operating condition\\
Switching condition & Switch-node, command, and current waveforms where available\\
Thermal behavior & Surface-temperature markers and their locations at recorded loads\\
\bottomrule
\end{tabularx}
\end{table}

The measurements were obtained during supervised prototype tests. USB-PD interoperability,
standardized safety and electromagnetic-emissions qualification, and a full
soft-switching map over input and load require separate test programs.
